\documentclass[10pt,a4paper]{amsart}
\usepackage[margin=19mm]{geometry}
\usepackage{amsmath,amssymb,mathtools}
\usepackage[scr=rsfs,cal=euler]{mathalfa}
\usepackage{xfrac}
\usepackage[unicode,hidelinks,bookmarks=false]{hyperref}
\newcommand{\ie}{i.e.\ }
\newcommand{\cf}{cf.\ }
\newcommand{\resp}{respectively\ }
\newcommand{\Subsection}{\subsection}
\providecommand{\nobreakdash}{\nobreak\hbox{-}}
\setlength{\emergencystretch}{2em}
\multlinegap0pt
\usepackage{amssymb}
\usepackage{enumerate}
\usepackage{tikz}
\usepackage{csquotes}
\usepackage{xcolor}
\usepackage[shortlabels]{enumitem}
\usepackage{float}
\newtheorem{proposition}{Proposition}
\newtheorem{theorem}{Theorem}
\usepackage{cleveref}
\crefname{proposition}{Proposition}{Propositions}
\crefname{theorem}{Theorem}{Theorems}
\newcommand{\Con}{{\mathrm{Con}}}
\def\F{\mathcal{F}}
\newcommand{\Supp}{\mathrm{Supp}}
\title{The complex of $r$-co-connected subgraphs, chordality and Fr\"oberg's theorem}
\author{Priyavrat Deshpande, Amit Roy, Rutuja Sawant}
\date{Source: arXiv:2510.25710v2 (CC BY 4.0)}
\begin{document}
\maketitle

\noindent\emph{Source and licence.} The complex of $r$-co-connected subgraphs, chordality and Fr\"oberg's theorem. Version arXiv:2510.25710v2 (CC BY 4.0), distributed by arXiv under the Creative Commons Attribution 4.0 International licence, http://creativecommons.org/licenses/by/4.0/, as recorded in the arXiv licence metadata for this exact version and shown on the abstract page https://arxiv.org/abs/2510.25710v2. Full licence notice: \texttt{licenses/arxiv-2510.25710-CC-BY-4.0.txt}.\par\smallskip

\noindent\textit{Editorial scope.} Counted unit: the article's proposition VD (source0.tex lines 1885--1887). The article's complete original proof of it is delivered with it (source0.tex lines 1888--1942, one \texttt{proof} environment of 55 lines). No mathematics is written, repaired or re-derived: the statement and the proof are the article's own lines, in the article's own notation, and no retained line is substituted or re-flowed.  Retained uncounted as the source-local context the proof needs: lines 827--833; lines 1465--1474; lines 1819--1824.  Nothing was omitted from any retained span, so the package carries no omission record.  No retained line contains a citation, so the wrapper carries no bibliography and no bibliography entry.  No retained line contains a figure environment, an image-inclusion command or drawing code, so no image is delivered and the package declares no asset dependency.  Editorial formatting only, in addition: an AMS \texttt{amsart} preamble that restates the article's own declarations and macros used by the retained lines, namely the environments \texttt{proposition}, \texttt{theorem} on separate counters, together with the standard packages those lines need.  The retained environments print in this document as Proposition 1, Theorem 1 in order of appearance, whereas the article prints the same items with the numbering of its own full document.  The complete native source of the article as distributed in the arXiv e-print for this exact version is bundled unchanged as \texttt{sources/188065/source0.tex}.
\begin{proposition}\label{Eq. Condition for a shedding vertex}
Let $G$ be a graph, $r$ a fixed positive integer, and $x \in V(\Sigma_r(A,G))$. The following conditions are equivalent:
\begin{enumerate}
	\item[(i)] $x$ is a shedding vertex of $\Sigma_r(A,G)$.		
	\item[(ii)] For every $F \in \Supp_r(A,G)$ with $x \notin F$, there exists a vertex $y \in F$ such that $(F \setminus \{y\}) \cup \{x\} \in \mathrm{Supp}_r(A,G)$.
\end{enumerate}
\end{proposition}

	\begin{theorem}\label{cycle vd theorem} 
		Let $n \geq 3$ and $r\geq 2$, and let $C_n$ be the cycle graph on $n$ vertices. The following statements are equivalent:
		\begin{enumerate}        
			\item $\Sigma_r(C_n)$ is vertex decomposable.
			\item $\Sigma_r(C_n)$ is shellable.
			\item $\Sigma_r(C_n)$ is Cohen-Macaulay.
			\item $n \leq r+2$.
			\item The clutter $\Con_r(C_n)$ is co-chordal.
		\end{enumerate}
	\end{theorem}

\begin{equation}\label{eq3146}
	\begin{split}
		V(G) &= \{x_1, x_2, \ldots, x_r, x_{r+1}, x_{r+2}\} \cup T, \quad \text{with} \quad T = \left\{ x_{r+3}, x_{r+4}, \ldots,  x_{(r+2)+\left\lfloor \tfrac{r+3}{2} \right\rfloor} \right\}, \\
		E(G) &= \left\{\{x_{r+2}, x_1\}, \{x_i, x_{i+1}\} \mid i \in [r+1] \right\}\cup \left\{ \left\{x_j, x_{r+2 + \tfrac{j+1}{2}} \right\} \mid j \in [r+2], j\text{ odd integer} \right\}.
		\end{split}
\end{equation}

	\begin{proposition}\label{shell but not VD}
		Let $G$ denote the graph with the vertex and circuit set described in \Cref{eq3146} and $r \ge 3$. Then $G$ is an $r$-gap-free unicyclic graph such that the complex $\Sigma_{r}(G)$ is shellable but not vertex decomposable.
	\end{proposition}

	\begin{proof}
For each $r\ge 8$, one can observe that $(r+2)+\lfloor\frac{r+3}{2}\rfloor<2r$. Thus, analyzing the remaining $3\le r\le 7$ cases individually, we see that $G$ is $r$-gap-free. 
	
	Our aim now is to show that $\Sigma_r(G)$ is not vertex decomposable. Indeed, for any $v\in V(G)$, we have $|N_G[v]|\le 3$. Thus, $E(\Con_r(G\setminus N_G[v])) \neq \emptyset$, and hence, by \Cref{Eq. Condition for a shedding vertex}, $\Sigma_r(G)$ has no shedding vertex. Consequently, $\Sigma_r(G)$ is not vertex decomposable.
    
		To show that $\Sigma_{r}(G)$ is shellable, we need to find a shelling order on the set of facets of $\Sigma_{r}(G)$, and we do this by defining first a special total order on $\mathcal P(T)$, the power set of $T$. Note that the elements $\{i\mid x_i\in T\}$ have natural order induced from the set of natural numbers, which we denote by $<_{\mathbb N}$. Now, we endow $\mathcal P(T)$ with a special order $\lessdot$ which is described in the following steps. 
		 
		 Let $I,J\in \mathcal P(T)$, where $I=\{x_{i_1}\ldots,x_{i_n}\}$ and $J=\{x_{j_1},\ldots,x_{j_m}\}$. Let $i=\max_{<_{\mathbb N}}\{i_1,\ldots,i_n\}$ and $j=\max_{<_{\mathbb N}}\{j_1,\ldots,j_m\}$ .
		 
		 \begin{enumerate}
		 	\item[{\bf Step 0:}] Declare $\emptyset\lessdot I$ for any $\emptyset\neq I\in \mathcal P(T)$.
		 	
		 	\item[{\bf Step 1:}] If $i<_{\mathbb N}j$, then declare $I\lessdot J$.
		 	
		 	\item[{\bf Step 2:}] If $i=j$, and $|I|<_{\mathbb N} |J|$, then declare $I\lessdot J$.
		 	
		 	\item[{\bf Step 3:}] If $i=j$ and $|I|=|J|$, then consider $I'=I\setminus\{x_i\}$ and $J'=J\setminus\{x_j\}$, and go back to Step 1. If $I'\lessdot J'$, then declare $I\lessdot J$.
		 \end{enumerate}
		As an example, if we consider the set $T=\{x_1,x_2,x_3,x_4\}$, then the elements of $\mathcal P(T)$ are ordered as follows:
		 \begin{align*}
		 \emptyset\lessdot\{x_1\}&\lessdot\{x_2\}\lessdot\{x_1,x_2\}\lessdot\{x_3\}\lessdot\{x_1,x_3\}\lessdot\{x_2,x_3\}\lessdot\{x_1,x_2,x_3\}\lessdot\{x_4\}\lessdot\{x_1,x_4\}\lessdot\\
		 &\{x_2,x_4\}\lessdot\{x_3,x_4\}\lessdot\{x_1,x_2,x_4\}\lessdot\{x_1,x_3,x_4\}\lessdot\{x_2,x_3,x_4\}\lessdot\{x_1,x_2,x_3,x_4\}.
		 \end{align*}
		 It is easy to see that $\lessdot$ is a total order on the elements of $\mathcal P(T)$. We now proceed to provide a shelling order $<$ on $\mathcal F(\Sigma_{r}(G))$ based on the ordering $(\mathcal P(T),\lessdot)$.
		 
		 Let $F\in \mathcal F(\Sigma_{r}(G))$. Then $F^c\in E(\Con_r(G))$. For each $I\in\mathcal P(T)$, we define $\mathcal F_I=\{F\in \mathcal{F}(\Sigma_{r}(G))\mid F^c\cap T=I\}$. Observe that $\mathcal{F}(\Sigma_{r}(G))=\sqcup_{I\in\mathcal P(T)}\mathcal F_I$, and $\mathcal F_{\emptyset}=\mathcal F(\Sigma_{r}(C_{r+2}))$, where $C_{r+2}$ is the induced subgraph of $G$ on the vertex set $\{x_1,\ldots,x_{r+2}\}$. By \Cref{cycle vd theorem}, the facets of $\Sigma_{r}(C_{r+2})$ has a shelling order. Based on these observations, we define the following:   
		 \begin{enumerate}
		 	\item[$\bullet$] The facets in $\mathcal F_{\emptyset}$ are given a shelling order.
		 	
		 	\item[$\bullet$] For each fixed $\emptyset\neq I\in\mathcal P(T) $, the facets in $\F_I$ are given any arbitrarily chosen but fixed total order.
		 	
		 	\item[$\bullet$] For any two distinct $I,J\in \mathcal P(T)$ with $F_1\in \F_I$ and $F_2\in \F_J$, we have $F_1<F_2$ if and only if $I\lessdot J$.
		 \end{enumerate}
		 
		 \noindent
		 {\bf Claim}: $(\mathcal F(\Sigma_{r}(G)), <)$ is a shelling order.
		 
		 \vspace{.1cm}
		 \noindent
		 {\it Proof of Claim}: By construction, if $F\in \F_{\emptyset}$ and $F$ is not the first facet in the given ordering, then for any $F'\in \F_{\emptyset}$ with $F'<F$, there exists some $F''\in \F_{\emptyset}$ with $F''<F$ such that $F'\cap F\subseteq F''\cap F$ and $\dim(F''\cap F)=\dim F-1$.
		 
		 Now, let $F\in \F_I$ for some $\emptyset\neq I\in\mathcal P(T)$. Take any $F_1\in \F(\Sigma_{r}(G))$ with $F_1<F$. Let $P:=F\cap C_{r+2}$ be the path with two end points $u,v\in V(C_{r+2})$. If $u\notin F_1$, then consider $F_u=(F\cup\{x\})\setminus\{u\}$, where $x\in I$. Observe that
		  \[F_u^c=(F^c\setminus\{x\})\cup\{u\}\in E(\Con_r(G)),\]
		   since $F^c\in E(\Con_r(G))$. Thus, $F_u\in \F(\Sigma_{r}(G))$ with $F_u<F$, $\dim(F_u\cap F)=\dim(F)-1$, and $F_1\cap F\subseteq F_u\cap F$. A similar argument can be given in case $v\notin F_1$.
		   
		   Next, we consider the case when $u,v\in F_1$. Then $F_1\cap F\supseteq P$ since $|V(G)|<2r+2$ for each $r\ge 4$, and the $r=3$ case can be easily checked separately. Let $Q$ be the path on the vertex set $V(C_{r+2}\setminus P)$, and define
		   \[
		   K=\{x_i\in T\mid \{x_i,z\}\in E(G)\text{ for some }z\in Q\}.
		   \]
		   Since $u,v\in F_1$ and $F_1^c\in E(\Con_r(G))$, we have $F_1\in \F_J$ for some $\emptyset\neq J\in \mathcal P(T)$ with $J\lessdot I$ and $J\neq I$. Moreover, $I,J\subseteq K$ and if $t=\max_{<_{\mathbb N}}\{i\mid x_i\in I\}$, then $J\subseteq \{x_i\in T\mid i\le t\}$. Observe that if $J\subseteq I$, then by our choice of ordering, $J=I$, since $u,v\in F_1$, which is a contradiction. Thus, $J\setminus I\neq\emptyset$, and hence, there exists $x_p\in J\setminus I$ such that $p<t$. In this case, if we take $\widehat F=(F\cup\{x_t\})\setminus\{x_p\}$, then 
		   \[
		   \widehat{F}^c=(F^c\setminus x_t)\cup\{x_p\}\in E(\Con_r(G)).
		   \]
		   Furthermore, $\widehat F\in \F(\Sigma_{r}(G))$ with $\widehat F<F$, $\dim(\widehat F\cap F)=\dim(F)-1$, and $F_1\cap F\subseteq \widehat F\cap F$. This completes the proof of the claim and subsequently, the proposition.
	\end{proof}

\end{document}
